\item \textbf{Día óptico versus día geométrico.}
La duración de un día óptico se define como el intervalo de tiempo entre el instante en que la parte superior del Sol naciente es apenas visible sobre el horizonte y el instante en que la parte superior del Sol apenas desaparece bajo el plano horizontal debido a la refracción atmosférica. La duración del día geométrico no incluye este efecto.
\begin{enumerate}
    \item Explique cuál es más largo, ¿un día óptico o un día geométrico?
    \item Encuentre la diferencia entre estos dos intervalos de tiempo. Modele la atmósfera de la Tierra como uniforme, con índice de refracción $1.000293$, una superficie superior definida con precisión y $8614\text{ m}$ de profundidad. Suponga que el observador está en el ecuador de la Tierra de modo que la trayectoria aparente del Sol es perpendicular al horizonte.
\end{enumerate}